%----------------------------------------------------------------------------
%  Saint Vincent Beamer theme -- demo deck
%
%  Build:  latexmk -pdf Main.tex
%          (runs bibtex automatically; needs two passes for the TOC)
%
%  Theme options:
%    \usetheme{SaintVincent}                    default (light)
%    \usetheme[night]{SaintVincent}             dark colour scheme
%    \usetheme[nosectionslides]{SaintVincent}   no divider slide per section
%    \usetheme[hidememo]{SaintVincent}          hide \memo / \NB annotations
%----------------------------------------------------------------------------
\documentclass[aspectratio=169]{beamer}

\usetheme{SaintVincent}

%----------------------------------------------------------------------------
%  NOTES
%----------------------------------------------------------------------------
%\setbeameroption{hide notes}                        % slides only
%\setbeameroption{show only notes}                   % notes only
%\setbeameroption{show notes on second screen=right} % both

%----------------------------------------------------------------------------
%  TITLE
%----------------------------------------------------------------------------
\title{Saint Vincent Theme}
\subtitle{A clean Beamer theme for economics}
\author{Gianluca Tisci}
\institute{Alma Mater Studiorum --- University of Bologna\\
           Bachelor in Economics and Finance}
\date{\today}

\renewcommand{\memouser}{Gianluca}

\begin{document}

\begin{frame}[plain, noframenumbering]
  \titlepage
\end{frame}

\begin{frame}{Overview}
  \tableofcontents
\end{frame}

%------------------------------------------------
\section{Theme basics}
%------------------------------------------------

\begin{frame}{Bullet points}{With a subtitle}
  \begin{itemize}
    \item First point
    \item Second point
      \begin{itemize}
        \item A nested point
      \end{itemize}
    \item Third point
  \end{itemize}
  \note[item]{Speaker notes still work as usual.}
\end{frame}

\begin{frame}{Blocks}
  Some text will be \alert{highlighted} because it matters. Do not overdo it.

  \begin{block}{Block}
    Carries the theme accent.
  \end{block}

  \begin{alertblock}{Alert block}
    The one place a second colour is warranted --- an alert that looks like
    everything else is not an alert.
  \end{alertblock}

  \begin{examples}
    Neutral, rather than introducing a third hue.
  \end{examples}
\end{frame}

\begin{frame}{Two columns}
  \begin{columns}[t]
    \column{.45\textwidth}
      \textbf{Heading}
      \begin{enumerate}
        \item Statement
        \item Explanation
        \item Example
      \end{enumerate}
    \column{.45\textwidth}
      The theme loads no font package, so slides use Beamer's own default
      family. Everything is styled by size, weight and colour only.
  \end{columns}
\end{frame}

%------------------------------------------------
\section{Shared features}
%------------------------------------------------

\begin{frame}{Pencil marks and results}
  These come from \texttt{almacore}, so they behave identically here, in the
  lecture-notes class and in the problem-set template.

  \medskip
  Underline a \lapis{key phrase}, or box \lapisbox{a definition}.

  \lapisboxeq{\max_{x}\; u(x) \quad \text{s.t.} \quad p \cdot x \le w}

  \result{$x^{\ast} = w / 2p$}
\end{frame}

\begin{frame}[fragile]{Call-outs, code and annotations}
  \begin{notebox}
    \boxentry{Why one core}
    The colours, pencil marks and call-outs live in a single class-agnostic
    package, so the slide theme and the notes class cannot drift apart.
  \end{notebox}

  \begin{code}
regress price mpg weight, robust
  \end{code}

  A draft note: \memo{tighten this slide} and \NB{add citation}.
\end{frame}

\begin{frame}{Table}
  \begin{table}
    \centering
    \begin{tabular}{lll}
      \toprule
      \textbf{Treatment} & \textbf{Response 1} & \textbf{Response 2} \\
      \midrule
      Treatment 1 & 0.0003262 & 0.562 \\
      Treatment 2 & 0.0015681 & 0.910 \\
      Treatment 3 & 0.0009271 & 0.296 \\
      \bottomrule
    \end{tabular}
    \caption{Table caption}
  \end{table}
\end{frame}

\begin{frame}{Figure}
  \begin{figure}
    \centering
    \begin{tikzpicture}
      \begin{axis}[econaxis, xlabel={$x_1$}, ylabel={$x_2$},
                   xmin=0, xmax=3.2, ymin=0, ymax=5.6,
                   width=8cm, height=5cm,
                   restrict y to domain=0:6]
        \addplot[almacurve, domain=0:3, samples=100] {-x^2 + 4};
        \addplot[almacurve, domain=0:3, samples=100] {-x^2 + 5};
        \addplot[almaconstraint, domain=0:3] {4.5 - 1.5*x};
      \end{axis}
    \end{tikzpicture}
    \caption{Indifference curves and a budget line.}
  \end{figure}
\end{frame}

\begin{frame}{Theorem}
  \begin{theorem}[Mass--energy equivalence]
    $E = mc^2$
  \end{theorem}
\end{frame}

\begin{frame}[fragile]{Citation}
  An example of the \verb|\cite| command:

  \medskip
  This statement requires a citation \cite{p1}.
\end{frame}

\begin{frame}{References}
  \footnotesize
  % \bibliographystyle must come before \bibliography, and the argument takes
  % no .bib extension -- the original had both the wrong way round, which is
  % why the citation never resolved.
  \bibliographystyle{apalike}
  \bibliography{reference}
\end{frame}

\begin{frame}[plain, noframenumbering]
  \vfill
  \centering{\Huge\bfseries\color{svaccent} Thank you}
  \vfill
\end{frame}

\end{document}
